% Standalone source-record documentation; no research attempts.
\documentclass[11pt]{article}
\usepackage[margin=25mm]{geometry}
\usepackage{fontspec,unicode-math}
\setmainfont{Latin Modern Roman}
\setmathfont{Latin Modern Math}
\usepackage{amsmath,amssymb,mathtools,xurl,hyperref}
\hypersetup{colorlinks=true,urlcolor=blue,linkcolor=blue}
\setcounter{secnumdepth}{0}
\setlength{\parindent}{0pt}
\setlength{\parskip}{5pt}
\setlength{\emergencystretch}{3em}
\newcommand{\sref}[2]{\hyperlink{source:#1:#2}{[S#2]}}
\newcommand{\cataloguescope}[1]{\paragraph{Source record.}#1}
\newcommand{\floor}[1]{\left\lfloor#1\right\rfloor}
\begin{document}
\section{Multiplicity Obstructions in Geometric Complexity Theory}\label{statement}

\subsection{Definitions and mathematical statement}

\paragraph{Shared notation.}Standard number systems are $\mathbb N,\mathbb Z,\mathbb Q,\mathbb R,\mathbb C$. Specialized notation and index ranges follow the cited source and the definitions below.


In the usual GCT setting, the general linear group $G=\operatorname{GL}(V)$ acts on a space of polynomials, with orbit $Gp$ and Zariski closure $\overline{Gp}$. Its coordinate ring is graded by polynomial degree. For an irreducible polynomial representation indexed by a partition $\lambda$, $\operatorname{mult}_\lambda\mathbb C[\overline{Gp}]_d$ is its multiplicity in the degree-$d$ part. Inclusion $\overline{Gp}\subseteq\overline{Gq}$ induces a surjection of coordinate rings in the reverse direction, so a larger multiplicity on the left obstructs such inclusion.
\paragraph{Record scope and qualifications.}
The catalogue title is retained as metadata. Named certificates, contracts, bounds or frameworks in that title are not asserted as results of the original source extract.
This is a family of multiplicity-obstruction questions. The supplied record does not specify $G$, $p$, $q$, or the size parameters; an answer must identify the particular geometric-complexity-theory family to which it applies.

Source-record limitation: The acting group and source polynomials p,q whose orbit closures are to be separated are not supplied; the core obstruction instance is undefined.
\paragraph{Statement recorded in the supplied source.}
For the polynomial orbit closures relevant to geometric complexity theory, can one find a partition $\lambda$ and degree $d$ such that
$$
\operatorname{mult}_\lambda\mathbb C[\overline{Gp}]_d>\operatorname{mult}_\lambda\mathbb C[\overline{Gq}]_d,
$$
thereby obtaining a multiplicity obstruction separating the orbit closures? \sref{statement}{2}
\paragraph{Source status.}
The supplied catalogue labels this source record open; that metadata was checked against the supplied snapshot on 2026-10-11. The label is preserved for provenance and is not an independent verification of present open status or mathematical correctness. The catalogue literature assessment is dated 2026-08-22.
\subsection{Short English statement}
Separate polynomial orbit closures by finding an irreducible representation occurring more often in one coordinate ring than in the other.
\subsection{Sources}
\begin{itemize}
\item[\hypertarget{source:statement:1}{S1}] Ulam AI and the UnsolvedMath Contributors, UnsolvedMath: A Curated Collection of Open Mathematics Problems, catalogue entry. Catalogue curation and metadata: CC BY 4.0. The source extract is retained; definitions, English summaries and qualifications are editorial additions.
\url{https://huggingface.co/datasets/ulamai/UnsolvedMath}.
\item[\hypertarget{source:statement:2}{S2}] Mini-Workshop: Kronecker, Plethysm, and Sylow Branching Coefficients and their Applications to Complexity Theory (2021), Oberwolfach Reports, DOI 10.4171/owr/2020/9, pp. 517--568
\url{https://doi.org/10.4171/owr/2020/9}.
\item[\hypertarget{source:statement:3}{S3}] Mulmuley and Sohoni, Geometric Complexity Theory I
\url{https://doi.org/10.1137/S009753970038715X}.
\item[\hypertarget{source:statement:4}{S4}] Ikenmeyer and Panova, Geometric complexity theory and matrix powering
\url{https://doi.org/10.1016/j.difgeo.2017.07.001}.
\item Source-specific Oberwolfach Reports licences, including CC BY-SA 4.0 where applicable, are retained; see THIRD\_PARTY\_NOTICES.md.
\end{itemize}
\paragraph{Preparation and review.}
The original source and UnsolvedMath contributors retain attribution for the source extract. Codex (GPT-6; exact serving revision not exposed) assisted with definitions, English summaries, scope notes and formatting. No new human mathematical authorship or independently verified solution is asserted. Catalogue identity checked against the complete supplied registry and current repository. Distinct upstream records, including repeated questions and broad or damaged extracts, are retained for complete coverage of the source CS category. A repeated record is not a new independent conjecture. Complete mathematical-equivalence and current-literature review has not been performed. Preparation and final review follow contribution prompts C and D, with source gaps disclosed. Structural and rendering checks do not establish a new theorem or open status.
\label{end:statement}
\end{document}
